% MAS1004 Assignment 1: long report.
%
% You fill this in with your agent, in your repository. AGENTS.md tells the
% agent how. Red text in the PDF is what still has to be replaced.
%
% - Pictures go in figures/. Copy them from results/, which is not in git.
% - Outputs go in outputs/. run.py saves what it prints in
%   results/<tag>_output.txt, and PLAN.md saves the other outputs in results/
%   with `| tee results/...`. Copy them here. \outputfile{outputs/check.txt}
%   prints a file exactly as it is. Every number in this report has to come
%   from such an output.
%
% To make the PDF: compress this folder into a zip, and on Overleaf choose
% New Project, Upload Project. Set Menu, Compiler to XeLaTeX, and Recompile.

\documentclass[11pt]{article}
\usepackage{mas1004}

\reporttitle{Assignment 1: Long report}
\studentname{Nordin Geels}
\studentid{G20260447}

\begin{document}
\makeheader

% ---------------------------------------------------------------------------
\section*{Problem 1. Categories and images}

\subsection*{Why these categories}
My classifier tells apart four sports by a photo of the match and its playing
area: professional soccer, rugby, ice hockey and basketball (pitch, pitch, rink
and court). I chose them because each one means something to me: I play rugby
myself and played at the World Cup with my team, I go to every PSV Eindhoven home
game, I saw ice hockey live in America, and I grew up with basketball. I chose
professional matches only.

None of the four is an ImageNet category: \texttt{imagenet\_classes.txt} has no
field, pitch, rink, court or stadium at all. It does have the balls (soccer ball,
rugby ball, basketball and others), which could be a shortcut when a ball is
large in the picture.

The search phrases became the folder names, so the classes are called
\texttt{professional\_soccer\_match\_stadium},
\texttt{professional\_rugby\_match\_stadium},
\texttt{professional\_ice\_hockey\_game\_arena} and
\texttt{professional\_basketball\_game\_arena}. The classes are the four sports,
recognised by their playing area; the names come from the search phrases.

\subsection*{What I expected the model to use}
From my notes of 7 October, written before training: mostly the colour, and the
difference in how you score points: rugby has long poles, soccer has a goal net
that is quite big compared to ice hockey, and basketball has the basket.

\subsection*{Images per category}
The downloader kept 150 images per category (the \texttt{downloaded} column below,
from \texttt{clean.py count}); the other columns are explained in Problem 4.

\outputfile{outputs/clean_count.txt}

% ---------------------------------------------------------------------------
\section*{Problem 2. Preparing an image}

\sidebyside{figures/prepare_original.jpg}{one of my own photos (IMG\_9681, soccer), 1078 x 727 pixels}
           {figures/prepare_result.png}{what prepare\_image made of it (224 x 224, shown with the colour scaling undone)}

The photo is wider than it is tall, so prepare\_image first makes the shorter side
256 pixels and then keeps only the 224 x 224 square in the middle. Everything
outside that square is gone: on the left part of the group of players at the
corner of the box, on the right the people walking along the side of the pitch,
and a thin strip at the top and bottom.

\sidebyside{figures/crop_cut_IMG_9691.jpg}{IMG\_9691 (basketball, 1290 x 1600): the yellow square is what prepare\_image keeps}
           {figures/crop_cut_IMG_9691_result.png}{what prepare\_image made of it}

In this tall phone photo the crop cuts off the bottom of the basketball court,
the thing I care about, while the big screen hanging from the roof stays almost
completely in. So the model sees more scoreboard than court. This is one of the
photos the model got wrong (Problem 5).

% ---------------------------------------------------------------------------
\section*{Problem 3. First training runs}

\begin{center}\footnotesize
\begin{tabular}{llllrrrr}
\toprule
tag & start & trained & epochs & lr & trainable parameters & train acc. & test acc. \\
\midrule
run & ImageNet & all layers & 10 & 0.0001 & 11,178,564 & 100.0\% & 89.2\% \\
scratch & random & all layers & 10 & 0.0001 & 11,178,564 & 100.0\% & 77.5\% \\
frozen & ImageNet & last layer & 10 & 0.001 & 2,052 & 92.3\% & 85.0\% \\
lr\_small & ImageNet & all layers & 10 & 1e-05 & 11,178,564 & 98.5\% & 80.0\% \\
\bottomrule
\end{tabular}
\end{center}

The fourth setting is my own choice: a smaller learning rate (\texttt{--lr 1e-5},
tag \texttt{lr\_small}), because I expected it to give the most accurate model.

\picturefile{figures/run_curves.png}

The loss went down. The train loss dropped from 0.842 in epoch 1 to 0.046 in
epoch 3 and kept creeping down to 0.004 in epoch 10. The test loss went down
from 0.532 to its lowest point, 0.322, in epoch 4, and then flattened out and
went up a little (0.355 in epoch 10). The train accuracy reached 1.000 in
epoch 3.

The train accuracy is 100.0\% and the test accuracy 89.2\%. The model got every
training image right, but not every test image: it has learned the training
images by heart, which works less well on images it has never seen.
That gap is the difference between remembering and generalising.

\begin{center}\picturefile[0.5\linewidth]{figures/run_confusion.png}\end{center}

The model mixes up rugby and soccer most: 8 rugby test images were called soccer.
That did not surprise me. Rugby and soccer are both played on grass, so when the
lines and the goal posts are not prominent in the picture, I already expected the
model to have trouble with them. Ice hockey and basketball did not surprise me
either when they got mixed up later: professional ice hockey and basketball arenas
have a similar layout, especially in America with the big screens in the middle,
and both have a light-coloured floor.

Starting from ImageNet gave 89.2\% on the test set, starting from random numbers
77.5\%, with the same images, code and number of epochs. I think the ImageNet
start is better because the network already knows how to see edges, shapes,
grass, crowds and lights from 1.2 million photos, while from random numbers it
has to learn all of that from only 480 training images. Both reached 100.0\%
on the training images, so the scratch model also learned the training images by
heart, but what it learned works less well on new images.

Training only the last layer (2,052 of the 11,178,564 parameters) reached 85.0\%
on the test set, close to the 89.2\% of training all layers, and it trained in 0.4
minutes instead of 1.5 minutes (see the outputs below). Most of the work is done by what the
network already learned on ImageNet.

My own setting, the smaller learning rate, did worse: 80.0\% on the test set.
Its test accuracy was still rising at epoch 10 (from 0.408 in epoch 1 to 0.800
in epoch 10, see its output below), so with such small steps 10 epochs were not
enough to finish learning. The first run with the normal learning rate was done
after 3 to 4 epochs.

\begin{center}
\picturefile[0.32\linewidth]{figures/scratch_curves.png}\hfill
\picturefile[0.32\linewidth]{figures/frozen_curves.png}\hfill
\picturefile[0.32\linewidth]{figures/lr_small_curves.png}\\
\small scratch \hfill frozen \hfill lr\_small
\end{center}

\subsection*{Output of the runs}
\outputfile{outputs/run_output.txt}
\outputfile{outputs/scratch_output.txt}
\outputfile{outputs/frozen_output.txt}
\outputfile{outputs/lr_small_output.txt}

% ---------------------------------------------------------------------------
\section*{Problem 4. Cleaning}

\subsection*{My rule}
An image belongs in a category only if it shows the playing area of that sport
(pitch, rink or court) and the playing area covers at least 25\% of the image.
Remove: (1) images of only the crowd or only the stands, without the playing
area; (2) stadiums or arenas photographed from the outside, because then you
cannot see the playing area, and a rugby stadium could just as well be a soccer
stadium; (3) images of any sport other than the four chosen (soccer, rugby, ice
hockey, basketball), for example American football.

Added later, during Step 14, after looking at the model's most confident
mistakes: (4) images that do not show a real playing area, for example an arcade
basketball machine. I decided to keep blurry images and images with watermarks.

\subsection*{Three kinds of junk}
\begin{center}
\picturefile[0.32\linewidth]{figures/junk_closeup_soccer_0057.jpg}\hfill
\picturefile[0.32\linewidth]{figures/junk_crowd_soccer_0072.jpg}\hfill
\picturefile[0.32\linewidth]{figures/junk_wrongsport_rugby_0061.jpg}
\end{center}
Left: soccer 0057, a close-up of two players with hardly any pitch (less than
25\% playing area; it is also a stock image with a watermark across it). Middle:
soccer 0072, fans in the stands, with the pitch only far away behind them (only
crowd or stands). Right: rugby 0061, an Australian rules football oval, not one
of my four sports.

\subsection*{What I removed}
\outputfile{outputs/clean_count.txt}

\outputfile{outputs/suspects_summary.txt}

The suspect list (run before I removed anything) listed 46 pairs of near copies.
From every pair one image is now removed: most with the reason
\textit{near copy of another image}, some because I had already removed one of
the two for another reason. In basketball there was a group of very similar
iStock photos; I kept four of them that are different moments from the same photo
shoot (same arena, same angle, same watermark), because the tool did not call
them copies of each other. Running \texttt{clean.py suspects} again after cleaning
found no near copies at all.

Of the first 20 suspects in each class I removed 4 (basketball), 7 (ice hockey),
9 (rugby) and 5 (soccer). Many suspects were fine pictures that were just
unusual, so I kept them.
The suspects did include junk I had missed when I looked through the images
myself, and I also found junk myself that the tool did not list. Looking myself
and using the suspect list both found things the other missed.

Two more images were removed later, at Step 14, after I saw them among the
model's most confident mistakes: basketball 0131 (a close-up of players, under
the 25\% rule) and basketball 0116 (an arcade basketball machine, rule 4). Both
were in the test set. Removing test images because the model got them wrong would
make the test score look better without the model improving, so I removed them
only because they break my rule.

\subsection*{Before and after cleaning}
\outputfile{outputs/compare.txt}

\begin{center}
\picturefile[0.49\linewidth]{figures/clean_curves.png}\hfill
\picturefile[0.4\linewidth]{figures/clean_confusion.png}
\end{center}

Before cleaning (\texttt{run}) the test accuracy was 89.2\% and the accuracy on my
new images 71.4\% (15/21). After cleaning (\texttt{clean}, same settings) it was
88.8\% and 81.0\% (17/21). The test accuracy went down a little, but cleaning
changed the test set itself: it is a different, smaller set of images without
the junk, so the two test numbers are not a fair comparison. My 21 new images are
exactly the same before and after, so that comparison is the fair one, and there
the model got two more images right.

The scores also move from run to run with the same data and settings. I trained
\texttt{clean} three times, and \texttt{run.py} and \texttt{check.py} printed:
89.7\% test and 85.7\% (18/21) new images the first time, 88.8\% and 81.0\%
(17/21) the second time (a re-run for Step 14), and 88.8\% and 81.0\% (17/21)
the third time (after 0116 and 0131 were also removed). Each run overwrote the
output files of the one before, so the files here are from the third run. The
score moves because training has random parts that the seed does not fix (the new
last layer starts from random numbers, and the training images are shuffled into
a new order every epoch), because the Apple Silicon GPU does not calculate in
exactly the same order every time, and because the test sets are small, so one
image more or less changes the percentage noticeably. In all three runs the clean
model got more of my new images right (18 or 17) than the model before cleaning
(15). I did not keep training until a higher number came out.

% ---------------------------------------------------------------------------
\section*{Problem 5. Images from a new source}

My 21 new images are my own photos and screenshots from my phone: the rugby
images are frames from my own rugby team's drone footage and the stream of our
match when we played at the World Cup, and photos from a rugby match I went to;
the basketball and ice hockey images are from games I went to in America; and the
soccer images are from my own phone gallery, for example from PSV home games.
None of them can be among my downloads: I took them or made the screenshots
myself, so they are not pictures the image search returned, and the overlap check
finds no near copies (closest 0.87; near copies are above 0.95). One rugby image
(IMG\_9678) is a screenshot of a stream of our own match, with ``PAUSED'' and a
broadcaster logo in it.

\outputfile{outputs/overlap.txt}

\begin{center}
\begin{tabular}{lrr}
\toprule
model & downloaded test set & my new images \\
\midrule
run (before cleaning) & 89.2\% & 71.4\% (15/21) \\
clean (after cleaning) & 88.8\% & 81.0\% (17/21) \\
\bottomrule
\end{tabular}
\end{center}

\picturefile{figures/own_photos_all.png}

All 5 rugby and all 6 soccer photos were right. The model got 4 of 5 ice hockey
photos and only 2 of 5 basketball photos right. On my new images it made only
four mistakes, shown below; as a fifth I add its most confident mistake on the
downloaded test set.

\picturefile{figures/own_photos_mistakes.png}

\begin{itemize}
\item IMG\_9695, basketball, called ice hockey (99\%): the floor is very pale,
almost white, like ice, and the whole hall is white.
\item IMG\_9697, ice hockey, called basketball (91\%): the ice is coloured dark red
by a light projection, so it no longer looks like white ice.
\item IMG\_9691, basketball, called ice hockey (62\%): a pale floor, and a big screen
hanging from the roof like in ice hockey arenas; the crop also cuts off the
bottom of the court (Problem 2).
\item IMG\_9694, basketball, called rugby (37\%): I am not sure, but maybe because
the photo has no crowd or stands, which in my analysis pushes the model towards
rugby, and the floor has a different colour from the darker, browner tones of
other basketball courts. The model itself was also unsure (rugby 37\%).
\item Rugby 0008 (downloaded test set), called soccer (100\%): a huge stadium full
of fans, with the rugby posts small in the picture.
\end{itemize}

What is different about my new images: they are phone photos taken from where I
sat or stood, often in portrait format, so the crop cuts off the top and bottom.
The lighting is real (arena lights, a red light show, night outside) instead of
the polished stock and AI-generated pictures that make up much of the downloads.
Some of my photos show a small hall or an outdoor court, while the downloads are
mostly big professional stadiums seen from above. Some are screenshots of a stream
or drone footage.

% ---------------------------------------------------------------------------
\section*{Problem 6. The demo}

Address: \url{https://geelsnordin-bot.github.io/mas1004-assignment1-nordin/}

The page shows the clean model; its self-test badge is green (the browser agrees
with Python, largest gap 0.0000), and it works on my phone over mobile data.

When I showed it a photo of food (rice with chicken), it still chose one of my
sports: soccer 79.1\%, rugby 15.2\%, ice hockey 4.3\%, basketball 1.4\%. It always
has to pick one of the four categories; it cannot say that a picture is none of
them.

% ---------------------------------------------------------------------------
\section*{Problem 7. What the model looks at}

I made every guess and chose every change myself; the agent only made the
changes I asked for. I asked in Dutch, so each request is given word for word in
Dutch, with an English translation. All model answers below are from the clean
model (the one on the demo page) and are printed together at the end of this
section. Every changed image is a copy saved in \texttt{data/changed/}; the
originals were never changed.

% --------------------------------------------------------------------------
\subsection*{Guess 1: the brightness of the floor}

\textbf{Guess.} I think the model uses the brightness of the floor: a pale or
white floor means ice hockey, a darker floor means basketball.

\textbf{Why.} IMG\_9695 (basketball, very pale floor) was called ice hockey with
99\%, and IMG\_9691 (basketball, pale floor) was called ice hockey, while IMG\_9693
and IMG\_9696, with a much darker floor, were recognised as basketball. And the
other way round: IMG\_9697, an ice hockey rink coloured dark red by a light
projection, was called basketball with 91\%.

\textbf{Change.} Recolour only the floor: pale basketball floors become a
brownish wood colour, and the red ice becomes plain white ice. If the guess is
right, the answers should flip. I also added three basketball test images that
the model called ice hockey or soccer (0143 with a grey floor, 0062 with a light
wooden floor, and 0064 with a green court).

My requests to the agent, word for word:
\begin{quote}
``Oké laten we dan kleur veranderen van de ice hockey foto zonder die lampen en
dus gewoon wit ijs en de basketbal velden een bruinigere kleur te geven de vloer
in plaats van zo bleek''\\
(Okay, then let's change the colour of the ice hockey photo without those lights,
so just white ice, and give the basketball courts a browner floor instead of so
pale.)

``Ja voeg die 2 ook maar toe aan de test plus 0064 en in plaats van dat groen dus
ook bruin maken''\\
(Yes, add those 2 to the test as well, plus 0064, and make that green brown too.)
\end{quote}

The floor was selected by an outline or by its colour and recoloured, keeping its
brightness pattern so the court lines stay visible; everything else is unchanged.

\sidebyside{figures/p7/new_images__professional_basketball_game_arena__IMG_9695.jpg}{IMG\_9695 (basketball), original: ice hockey 99\%, basketball 0\%}
           {figures/p7/changed__basketball_brown_floor__IMG_9695.jpg}{brown floor: ice hockey 86\%, basketball 11\%}
\sidebyside{figures/p7/new_images__professional_basketball_game_arena__IMG_9691.jpg}{IMG\_9691 (basketball), original: ice hockey 62\%, basketball 36\%}
           {figures/p7/changed__basketball_brown_floor__IMG_9691.jpg}{brown floor: basketball 74\%, ice hockey 25\%}
\sidebyside{figures/p7/clean__professional_basketball_game_arena__0143.jpg}{0143 (basketball), original: ice hockey 96\%, rugby 2\%}
           {figures/p7/changed__basketball_brown_floor__0143.jpg}{brown floor: ice hockey 44\%, basketball 34\%}
\sidebyside{figures/p7/clean__professional_basketball_game_arena__0062.jpg}{0062 (basketball), original: ice hockey 83\%, basketball 11\%}
           {figures/p7/changed__basketball_brown_floor__0062.jpg}{brown floor: basketball 77\%, ice hockey 19\%}
\sidebyside{figures/p7/clean__professional_basketball_game_arena__0064.jpg}{0064 (basketball), original: soccer 54\%, basketball 41\%}
           {figures/p7/changed__basketball_brown_floor__0064.jpg}{green parts brown: basketball 100\%, soccer 0\%}
\sidebyside{figures/p7/new_images__professional_ice_hockey_game_arena__IMG_9697.jpg}{IMG\_9697 (ice hockey), original: basketball 91\%, ice hockey 8\%}
           {figures/p7/changed__ice_hockey_white_ice__IMG_9697.jpg}{plain white ice: ice hockey 69\%, basketball 30\%}

\textbf{What it says.} Four of the six images flipped to the right sport, 0143
nearly did, and IMG\_9695 stayed ice hockey but became less sure. In IMG\_9697 the
white ice also removed the players and the projected logo, so more than the colour
changed there. My conclusion: the colour of the playing surface has a lot of
influence on the model.

% --------------------------------------------------------------------------
\subsection*{Guess 2: the brightness of the surroundings}

\textbf{Guess.} I think the model also uses how bright the surroundings are: a
very light place (a white hall, or outside under a light blue sky) makes it think
of ice hockey.

\textbf{Why.} IMG\_9695 (white hall) and 0143 (outside, light sky) were still
called ice hockey after their floor was made brown in guess 1.

\textbf{Change.} Starting from the brown-floor versions, make everything except
the floor darker and browner (test 2). Then, for IMG\_9695, make everything except
the floor black (test 3).

My requests to the agent, word for word:
\begin{quote}
``Op de bruine omgeving de setting donkerder maken met bijvoorbeeld tribunes met
met veel mensen die wat donkerder/bruiner zijn getint dan witte muren bijvoorbeeld
ik denk dat dit veel gaat helpen met de gok omdat dan het hoge wit gehalte weg is
uit de foto''\\
(On the brown version, make the setting darker, for example with stands full of
people that are tinted darker/browner than white walls; I think this will help a
lot because then the high amount of white is gone from the photo.)

``Conclusie is dat het helpt kan je het nog is testen door die witte muren gewoon
zwart te maken zodat er alleen gefocust wordt op het sport veld wat dus al bruiner
is?''\\
(The conclusion is that it helps. Can you test it again by simply making those
white walls black, so that only the sports floor, which is already browner, is
focused on?)

``[...] maar je hebt nu ook de basket er uit gedaan dat is niet de bedoeling
[...]''\\
([...] but you have now also taken out the basket, that was not the intention
[...])
\end{quote}

In test 2 the surroundings were made darker and browner by a colour tint; no
people or stands were added. In test 3 the agent first also blacked out the basket,
which I had not asked for, so it was done again with the basket kept.

\sidebyside{figures/p7/changed__basketball_brown_floor__IMG_9695.jpg}{IMG\_9695, brown floor: ice hockey 86\%, basketball 11\%}
           {figures/p7/changed__basketball_brown_floor_dark_surroundings__IMG_9695.jpg}{+ dark surroundings: ice hockey 85\%, basketball 10\%}
\sidebyside{figures/p7/changed__basketball_brown_floor__0143.jpg}{0143, brown floor: ice hockey 44\%, basketball 34\%}
           {figures/p7/changed__basketball_brown_floor_dark_surroundings__0143.jpg}{+ dark surroundings: basketball 80\%, soccer 10\%}
\sidebyside{figures/p7/changed__basketball_brown_floor__IMG_9695.jpg}{IMG\_9695, brown floor: ice hockey 86\%, basketball 11\%}
           {figures/p7/changed__basketball_brown_floor_black_surroundings__IMG_9695.jpg}{black, basket also removed: ice hockey 50\%, rugby 20\%}
\sidebyside{figures/p7/changed__basketball_brown_floor__IMG_9695.jpg}{IMG\_9695, brown floor: ice hockey 86\%, basketball 11\%}
           {figures/p7/changed__basketball_brown_floor_black_surroundings__IMG_9695_with_basket.jpg}{black, basket kept: ice hockey 44\%, basketball 29\%}

\textbf{What it says.} Darkening the surroundings helps: 0143 flipped to
basketball. But IMG\_9695 barely changed, and even with only the floor (and the
basket) left it stayed ice hockey, though much less sure. My conclusion: it does
not really have to do with the white walls.

% --------------------------------------------------------------------------
\subsection*{Guess 3: an empty court}

\textbf{Guess.} The court in IMG\_9695 looks empty, while almost all other
basketball photos look full, because a basketball court is small compared to the
fields of the other sports. Adding a ball and players might make the model say
basketball.

\textbf{Change.} On the brown-floor version of IMG\_9695, first add a ball, then
players (test 4). Then the other way round: remove the players from basketball
test images that the model got right (test 5).

My requests to the agent, word for word:
\begin{quote}
``[...] plus was misschien helpt is door een basketbal toe toevoegen op het veld
met wat spelers? Het veld is nu namelijk erg leeg terwijl bijna alle andere
foto’s met basketbal best vol lijken door het kleine veldje dat het is als je het
vergelijkt met de andere sporten''\\
([...] and what might help is adding a basketball on the court with some players?
The court is very empty now, while almost all other basketball photos look quite
full because of how small the court is compared with the other sports.)

``Test A en dan B''\\
(Test A and then B: first the ball, then the players.)

``Dat spelers op een veld dus kunnen helpen om het goed te herkennen we kunnen
misschien van een goed gegokte spelers weghalen om te kijken wat er dan gebeurt
met de score?''\\
(So players on a court can help to recognise it correctly. Maybe we can remove the
players from a correctly guessed one, to see what happens to the score?)

``Ja doe maar met testset en kies welke het minst nep lijkt of het model verkeerd
beïnvloed''\\
(Yes, use the test set, and choose the ones that look least fake or influence the
model wrongly the least.)
\end{quote}

The ball is drawn. The players were cut out of basketball image 0020 (a near copy
removed from the training data, so not a test image) as rectangles with soft
edges, so they also bring a bit of dark crowd and yellow floor from that photo.
For test 5 the agent chose 0002, 0090 and 0101, and painted the players out with
the colours to their left and right (crowd stays crowd, floor stays floor), which
leaves visible smears; the ball, referees and spectators stayed.

\sidebyside{figures/p7/changed__basketball_brown_floor__IMG_9695.jpg}{IMG\_9695, brown floor: ice hockey 86\%, basketball 11\%}
           {figures/p7/changed__basketball_brown_floor_ball__IMG_9695.jpg}{+ ball: ice hockey 86\%, basketball 11\%}
\sidebyside{figures/p7/changed__basketball_brown_floor__IMG_9695.jpg}{IMG\_9695, brown floor: ice hockey 86\%, basketball 11\%}
           {figures/p7/changed__basketball_brown_floor_players__IMG_9695.jpg}{+ players: basketball 58\%, ice hockey 40\%}
\sidebyside{figures/p7/changed__basketball_brown_floor__IMG_9695.jpg}{IMG\_9695, brown floor: ice hockey 86\%, basketball 11\%}
           {figures/p7/changed__basketball_brown_floor_ball_players__IMG_9695.jpg}{+ ball and players: basketball 51\%, ice hockey 46\%}
\sidebyside{figures/p7/clean__professional_basketball_game_arena__0002.jpg}{0002 (basketball), original: basketball 100\%, rugby 0\%}
           {figures/p7/changed__basketball_players_removed__0002.jpg}{players removed: basketball 96\%, ice hockey 2\%}
\sidebyside{figures/p7/clean__professional_basketball_game_arena__0090.jpg}{0090 (basketball), original: basketball 100\%, ice hockey 0\%}
           {figures/p7/changed__basketball_players_removed__0090.jpg}{players removed: basketball 92\%, soccer 5\%}
\sidebyside{figures/p7/clean__professional_basketball_game_arena__0101.jpg}{0101 (basketball), original: basketball 100\%, rugby 0\%}
           {figures/p7/changed__basketball_players_removed__0101.jpg}{players removed: basketball 100\%, soccer 0\%}

\textbf{What it says.} The ball did nothing, but with players added IMG\_9695 was
called basketball for the first time. My conclusion for test 4: players on a court
can help the model recognise it correctly. Removing players barely changed the
answer (all three stayed basketball). My conclusion for test 5: maybe adding
players is easier and disturbs the photo less than removing them; players that are
removed always stay somewhat visible in the photo.

% --------------------------------------------------------------------------
\subsection*{Guess 4: the rugby posts}

\textbf{Guess.} I think the model recognises rugby by the tall H-shaped posts:
when the posts are clearly visible it says rugby, and when they are small or not
prominent it says soccer.

\textbf{Why.} I compared the soccer and rugby mistakes with test images it got
right. In the correctly recognised rugby photos you clearly see the posts or
players close together. In the rugby photos it called soccer, like 0008, the posts
are not prominent. And soccer 0107 was called rugby, where tall thin poles above
the stadium look like posts.

\begin{center}\picturefile{figures/soccer_rugby_compare.png}\end{center}

\textbf{Change.} Zoom in on only the field in 0008 (posts should become more
prominent, so I expected rugby) and in 0107 (without the poles above the stadium,
so hopefully soccer). Then test it on more images: replace the goal in soccer
photos the model got right with rugby posts (test 6b).

My requests to the agent, word for word:
\begin{quote}
``Oké test met 0008 de palen door het in te zoomen op alleen het veld en test dat
ook bij 0107 alleen het veld want die dingen boven het stadion laten het lijken als
palen dus alleen op veld in zoomen en dan hopelijk gokt het soccer [...]''\\
(Okay, test the posts with 0008 by zooming in on only the field, and test that also
with 0107, only the field, because the things above the stadium make it look like
posts; so only zoom in on the field, and then hopefully it guesses soccer [...])

``Oké zoom verder in zodat je niet de tribune ziet en alleen het veld met de
palen''\\
(Okay, zoom in further so that you do not see the stands, only the field with the
posts.)

``Test 6 zoek een paar voetbal foto’s en voeg in plaats van het goal daar palen
toe. [...]''\\
(Test 6: find a few soccer photos and add posts there instead of the goal. [...])

``[...] Plus bij die rugby palen foto moet er wel alleen op het veld ingezoomd
worden aangezien veel publiek al snel een hoge voetbal score geeft''\\
([...] And for the rugby posts photos you have to zoom in on only the field,
because a lot of crowd quickly gives a high soccer score.)
\end{quote}

The first zoom of 0008 was wide, and the model only sees the middle square, where
the posts were at the edge, so a square crop with the posts in the middle was made
as well (I asked for that too: ``0008 opnieuw met palen in het midden'', 0008 again
with the posts in the middle), and then the close-up. Because the posts stand in front of the
stands, the close-up still shows some stands at the top. In test 6b the goal was
covered with the grass next to it and white H-shaped posts, about five times as
tall as the goal, were drawn in its place; in 0018 part of a second goal stayed
visible.

\sidebyside{figures/p7/clean__professional_rugby_match_stadium__0008.jpg}{0008 (rugby), original: soccer 100\%, rugby 0\%}
           {figures/p7/changed__zoom_field_only__0008_rugby.jpg}{zoomed in on the field: soccer 99\%, rugby 1\%}
\sidebyside{figures/p7/clean__professional_rugby_match_stadium__0008.jpg}{0008 (rugby), original: soccer 100\%, rugby 0\%}
           {figures/p7/changed__zoom_field_only__0008_rugby_posts_centred.jpg}{square, posts in the middle: soccer 95\%, rugby 5\%}
\sidebyside{figures/p7/clean__professional_rugby_match_stadium__0008.jpg}{0008 (rugby), original: soccer 100\%, rugby 0\%}
           {figures/p7/changed__zoom_field_only__0008_rugby_posts_close.jpg}{close-up of posts and field: rugby 99\%, soccer 1\%}
\sidebyside{figures/p7/clean__professional_soccer_match_stadium__0107.jpg}{0107 (soccer), original: rugby 77\%, soccer 19\%}
           {figures/p7/changed__zoom_field_only__0107_soccer.jpg}{zoomed in on the field: soccer 98\%, basketball 1\%}
\sidebyside{figures/p7/clean__professional_soccer_match_stadium__0002.jpg}{0002 (soccer), original: soccer 100\%, rugby 0\%}
           {figures/p7/changed__soccer_goal_to_rugby_posts__0002.jpg}{goal replaced by rugby posts: soccer 100\%, rugby 0\%}
\sidebyside{figures/p7/clean__professional_soccer_match_stadium__0002.jpg}{0002 (soccer), original: soccer 100\%, rugby 0\%}
           {figures/p7/changed__soccer_goal_to_rugby_posts_zoomed__0002.jpg}{posts, zoomed in on the field: rugby 89\%, soccer 10\%}
\sidebyside{figures/p7/clean__professional_soccer_match_stadium__0018.jpg}{0018 (soccer), original: soccer 100\%, basketball 0\%}
           {figures/p7/changed__soccer_goal_to_rugby_posts__0018.jpg}{goal replaced by rugby posts: soccer 100\%, rugby 0\%}
\sidebyside{figures/p7/clean__professional_soccer_match_stadium__0018.jpg}{0018 (soccer), original: soccer 100\%, basketball 0\%}
           {figures/p7/changed__soccer_goal_to_rugby_posts_zoomed__0018.jpg}{posts, zoomed in on the field: soccer 52\%, rugby 47\%}
\sidebyside{figures/p7/clean__professional_soccer_match_stadium__0055.jpg}{0055 (soccer), original: soccer 100\%, rugby 0\%}
           {figures/p7/changed__soccer_goal_to_rugby_posts__0055.jpg}{goal replaced by rugby posts: soccer 100\%, rugby 0\%}
\sidebyside{figures/p7/clean__professional_soccer_match_stadium__0055.jpg}{0055 (soccer), original: soccer 100\%, rugby 0\%}
           {figures/p7/changed__soccer_goal_to_rugby_posts_zoomed__0055.jpg}{posts, zoomed in on the field: rugby 72\%, soccer 26\%}

\textbf{What it says.} Without the poles above the stadium, 0107 became soccer.
For 0008 the answer only changed to rugby in the close-up where the posts were big.
In the soccer photos the drawn posts did nothing in the whole stadium view, where
they are only thin lines, but after zooming in on the field two of the three
became rugby and the third (0018, with the most stands behind the posts) nearly
did. Zooming in changes two things at once: the posts get bigger and the crowd gets
smaller. My conclusion: posts that are prominent in the picture influence the
rugby score, while crowd and stands can push towards a higher soccer score; this
shows in 6b, where the zoomed photo with the most crowd behind the prominent posts
(0018) has the lowest rugby score.

% --------------------------------------------------------------------------
\subsection*{Guess 5: players close together}

\textbf{Guess.} I think the model recognises rugby by players standing close
together in a group, like a lineout or a scrum.

\textbf{Why.} The same comparison as guess 4: in the rugby photos it got right you
often see players close together. And soccer 0137, with players spread out, was
called rugby.

\textbf{Change.} Remove the players from 0137, and remove only the group of players
standing close together from rugby test images the model got right (test 7b).

My requests to the agent, word for word:
\begin{quote}
``[...] en test de spelers dicht bij elkaar met 0137 door de spelers weg te
laten''\\
([...] and test the players close together with 0137 by leaving the players out.)

``[...] Test 7 doe dit bij andere rugby velden had namelijk bij de test geen invloed
[...]''\\
([...] Test 7: do this with other rugby fields, because in the test it had no
influence [...])

``Je hebt 7 b niet goed gedaan er zijn tribunes toegevoegd?????''\\
(You did 7b wrong, stands were added?????)
\end{quote}

The agent first misread ``test 7'' as the crowd test and added stands to rugby
photos (shown under guess 6); after my correction it removed the player groups:
the group in 0012, the scrum in 0028 and the lineout in 0059. Other players and
the referee stayed. The painted area leaves visible smears, very clearly in 0059.
I expected little effect.

\sidebyside{figures/p7/clean__professional_soccer_match_stadium__0137.jpg}{0137 (soccer), original: rugby 51\%, soccer 36\%}
           {figures/p7/changed__soccer_players_removed__0137.jpg}{players removed: rugby 42\%, soccer 33\%}
\sidebyside{figures/p7/clean__professional_rugby_match_stadium__0012.jpg}{0012 (rugby), original: rugby 71\%, soccer 28\%}
           {figures/p7/changed__rugby_player_group_removed__0012.jpg}{player group removed: rugby 84\%, soccer 15\%}
\sidebyside{figures/p7/clean__professional_rugby_match_stadium__0028.jpg}{0028 (rugby), original: rugby 100\%, soccer 0\%}
           {figures/p7/changed__rugby_player_group_removed__0028.jpg}{player group removed: rugby 70\%, soccer 29\%}
\sidebyside{figures/p7/clean__professional_rugby_match_stadium__0059.jpg}{0059 (rugby), original: rugby 94\%, basketball 5\%}
           {figures/p7/changed__rugby_player_group_removed__0059.jpg}{player group removed: rugby 86\%, basketball 10\%}

\textbf{What it says.} All of them stayed with the answer they had (0137 stayed
rugby, the three rugby photos stayed rugby, 0028 less sure). My conclusion:
removing the players standing close together does not necessarily make the model
think it is another sport.

% --------------------------------------------------------------------------
\subsection*{Guess 6: the amount of crowd}

\textbf{Guess.} A lot of crowd in the picture makes the model say soccer, little
crowd makes it say rugby.

\textbf{Why.} My rugby photos hardly show a crowd, while the soccer photos show a
lot, and the close-up of 0008, with less crowd, became rugby.

\textbf{Change.} Add more stands with crowd to 0137 (test 7), and remove the stands
and crowd from soccer photos the model got right (test 8).

My requests to the agent, word for word:
\begin{quote}
``Bij 0137 dan meer tribune/ publiek toevoegen''\\
(For 0137, add more stands/crowd.)

``[...] en test 8 haal bij voetbal foto’s de tribunes weg en geen publiek''\\
([...] and test 8: remove the stands from soccer photos, and no crowd.)
\end{quote}

For 0137 the stands at the top of the photo itself were repeated above it until the
picture was square. This looks repetitive, and because the picture became square
the model now sees the full width of the pitch, so two things changed at once. For
test 8 each soccer photo was cropped to the pitch only; in 0060 a little of the
stands is still visible.

\sidebyside{figures/p7/clean__professional_soccer_match_stadium__0137.jpg}{0137 (soccer), original: rugby 51\%, soccer 36\%}
           {figures/p7/changed__soccer_more_crowd__0137.jpg}{more stands added on top: soccer 82\%, rugby 13\%}
\sidebyside{figures/p7/clean__professional_soccer_match_stadium__0115.jpg}{0115 (soccer), original: soccer 100\%, rugby 0\%}
           {figures/p7/changed__soccer_no_stands__0115.jpg}{only the pitch, no stands: soccer 87\%, rugby 11\%}
\sidebyside{figures/p7/clean__professional_soccer_match_stadium__0060.jpg}{0060 (soccer), original: soccer 98\%, rugby 2\%}
           {figures/p7/changed__soccer_no_stands__0060.jpg}{only the pitch, no stands: soccer 68\%, rugby 31\%}
\sidebyside{figures/p7/clean__professional_soccer_match_stadium__0092.jpg}{0092 (soccer), original: soccer 100\%, rugby 0\%}
           {figures/p7/changed__soccer_no_stands__0092.jpg}{only the pitch, no stands: soccer 78\%, rugby 21\%}

When the agent misunderstood test 7b, it added stands with crowd to three rugby
photos in the same way. I did not ask for this, but the answers are useful for this
guess: all three stayed rugby (originals: 0012 rugby 71\%, 0028 rugby 100\%, 0063
rugby 100\%, see the outputs at the end of this section).

\sidebyside{figures/p7/clean__professional_rugby_match_stadium__0012.jpg}{0012 (rugby), original: rugby 71\%, soccer 28\%}
           {figures/p7/changed__rugby_more_crowd_misunderstood_7b__0012.jpg}{stands with crowd added: rugby 68\%, soccer 31\%}
\sidebyside{figures/p7/clean__professional_rugby_match_stadium__0028.jpg}{0028 (rugby), original: rugby 100\%, soccer 0\%}
           {figures/p7/changed__rugby_more_crowd_misunderstood_7b__0028.jpg}{stands with crowd added: rugby 99\%, basketball 1\%}
\sidebyside{figures/p7/clean__professional_rugby_match_stadium__0063.jpg}{0063 (rugby), original: rugby 100\%, soccer 0\%}
           {figures/p7/changed__rugby_more_crowd_misunderstood_7b__0063.jpg}{stands with crowd added: rugby 100\%, soccer 0\%}

\textbf{What it says.} With more stands, 0137 became soccer (one image, with two
changes at once). Removing the stands from soccer photos kept them soccer but made
the model less sure, moving them towards rugby. Adding stands to rugby photos did
not make them soccer. My conclusion: clearly visible posts lead to a higher rugby
score, while a lot of crowd and stands pushes the answer towards soccer.

% --------------------------------------------------------------------------
\subsection*{The model's answers for every image in this section}
\outputfile{outputs/problem7_predictions.txt}

% --------------------------------------------------------------------------
\subsection*{What I now think the model uses}
My tests suggest that the model relies heavily on the setting of a sport rather
than on the ball. Changing the colour of the playing surface was enough in most
images to shift the predictions towards basketball and ice hockey. For rugby,
prominent goal posts combined with few stands or spectators made the biggest
difference. Football worked the other way around: removing the stands around a
green grass pitch made the model less sure about football and moved it towards
rugby. Together, these results indicate that the classifier has mainly learned to
recognise the environment in which each sport is usually played, and that these
background cues can outweigh the action in the image itself.

I am fairly confident in this conclusion. The effects of the floor colour and the
goal posts (when they were prominent in the picture) showed up consistently across
several different images, which makes it unlikely that they were one-off results.
At the same time, some of my tests used only one to three images, so the evidence
for those specific changes is limited. A larger set of edited images would be
needed to confirm them.

% ---------------------------------------------------------------------------
\section*{Appendix. Output of check.py}
\outputfile{outputs/check.txt}

\end{document}
